% !TEX program = xelatex
% À compiler sur Overleaf avec le compilateur XeLaTeX
% (Menu -> Compiler -> XeLaTeX) pour avoir la police Verdana.

\documentclass[12pt,a4paper]{article}

\usepackage{fontspec}
% Verdana si disponible, sinon une police sans-serif proche
\IfFontExistsTF{Verdana}{\setmainfont{Verdana}}{\setmainfont{DejaVu Sans}}

\usepackage[french]{babel}
\usepackage[margin=1.4cm]{geometry}
\usepackage{amsmath}
\usepackage{setspace}
\usepackage{enumitem}
\usepackage{array}

% Interligne 1,5
\onehalfspacing

% Espacement réduit autour des égalités successives
\setlength{\abovedisplayskip}{1pt}
\setlength{\belowdisplayskip}{1pt}
\setlength{\abovedisplayshortskip}{0pt}
\setlength{\belowdisplayshortskip}{0pt}

% ----- Pas de numérotation des pages -----
\pagestyle{empty}

\setlength{\parindent}{0pt}

% ----- Commandes personnalisées -----
\newcounter{exercice}
\newcommand{\exercice}{%
  \refstepcounter{exercice}%
  \vspace{0.4em}
  \noindent\underline{\textbf{Exercice \arabic{exercice}}}\par
  \vspace{0em}
}

\setlist[enumerate,1]{label=\alph*),itemsep=0pt,topsep=1pt}
\setlist[itemize,1]{leftmargin=1.6em,itemsep=0pt,topsep=1pt}

\begin{document}

% ----- En-tête -----
\begin{center}
  \underline{\Large\textbf{Correction du contrôle séquence 13 version 5}}\\[0.3em]
 
\end{center}

% =====================================================================
\vspace{3cm}\exercice
\begin{enumerate}
  \item \textbf{Oui.} $45  = 5 \times 9$ 
  \newline 5 est donc un diviseur de 45.
  \item \textbf{Non.} $187 \div 6 \approx 31{,}17$
  \newline le quotient n'est pas un entier,
        donc 187 n'est pas un multiple de 6.
  \item \textbf{Non.}  $25 < 75$, donc 25 n'est pas un multiple de 75.
  \item \textbf{Non.} $54 \div 8 = 6{,}75$ 
  \newline le quotient n'est pas un entier,
        donc 8 n'est pas un diviseur de 54.
\end{enumerate}

% =====================================================================
\vspace{1cm}\exercice
\begin{enumerate}
  \item \textbf{2 et 3} (acceptés aussi : 5 ; 7 ; 11 ; \dots).
        Un nombre premier a exactement deux diviseurs : 1 et lui-même.
  \item \textbf{4 et 6} (acceptés aussi : 1 ; 8 ; 9 ; \dots).
      
  \item Dans 14 ; 22 ; 11 ; 27 ; 49 ; 58 ; 83 ; 121, les deux nombres premiers sont
        \textbf{11} et \textbf{83}.
       
\end{enumerate}

% =====================================================================
\vspace{1cm}\exercice
\textbf{a) } $21 = 3 \times 7$, 
donc les diviseurs de 21 sont $\boxed{1\ ;\ 21\ ;\ 3\ ;\ 7\ }$.

\textbf{b) } $28 = 2 \times 2 \times 7$,
donc les diviseurs de 28 sont
$\boxed{1\ ;\ 28\ ;\ 2\ ;\ 7\ ;\ 4\ ;\ 14\ }$.

% =====================================================================
\vspace{1cm}\exercice


\noindent\begin{minipage}[t]{0.48\textwidth}
\textbf{a)}
\begin{align*}
26 &= 2 \times 13
\end{align*}

\textbf{b)}
\begin{align*}
63 &= 3 \times 21\\
   &= 3 \times 3 \times 7
\end{align*}
\end{minipage}\hfill
\begin{minipage}[t]{0.48\textwidth}
\textbf{c)}
\begin{align*}
140 &= 2 \times 70\\
    &= 2 \times 2 \times 35\\
    &= 2 \times 2 \times 5 \times 7\\
\end{align*}

\textbf{d)}
\begin{align*}
182 &= 2 \times 91\\
    &= 2 \times 7 \times 13
\end{align*}
\end{minipage}

% =====================================================================
\exercice
\noindent\begin{minipage}[t]{0.31\textwidth}
\textbf{a)} 
\begin{align*}
\dfrac{6}{18} &= \dfrac{6 \times 1}{6 \times 3}\\
              &= \dfrac{1}{3}
\end{align*}
\end{minipage}\hfill
\begin{minipage}[t]{0.31\textwidth}
\textbf{b)} 
\begin{align*}
\dfrac{10}{15} &= \dfrac{5 \times 2}{5 \times 3}\\
               &= \dfrac{2}{3}
\end{align*}
\end{minipage}\hfill
\begin{minipage}[t]{0.31\textwidth}
\textbf{c)} 
\begin{align*}
\dfrac{75}{165} &= \dfrac{15 \times 5}{15 \times 11}\\
                &= \dfrac{5}{11}
\end{align*}
\end{minipage}

% =====================================================================
\vspace{1cm}\exercice
Le nombre de paquets doit diviser à la fois 135 et 105 : c'est un diviseur commun à
135 et 105. Pour en faire un maximum, on cherche le plus grand diviseur commun à 135 et 105.

Avec les décompositions en facteurs premiers :
$135 = 3 \times 3 \times 3 \times 5$ et $105 = 3 \times 5 \times 7$.
Le plus grand diviseur commun à 135 et 105 est le produit des facteurs premiers communs :
3 \times 5
                      = 15

Hulk et Batman peuvent donc faire au maximum \textbf{15 paquets},
chacun contenant  9 cookies et 7 tartelettes.

% =====================================================================
\vspace{1cm}\exercice
On utilise les critères de divisibilité :
\newline \textbf{par 2} : le chiffre des unités est 0, 2, 4, 6 ou 8 ;
\newline \textbf{par 3} : la somme des chiffres est divisible par 3
(sommes des chiffres : $31$ pour 671\,890 ; $54$ pour 666\,666\,666 ; $40$ pour 675\,985) ;
\newline \textbf{par 5} : le chiffre des unités est 0 ou 5 ;
\newline \textbf{par 10} : le chiffre des unités est 0.

\renewcommand{\arraystretch}{1.4}
\setlength{\tabcolsep}{2pt}
\begin{center}
\small
\begin{tabular}{|c|c|c|c|c|}
\hline
 & \textbf{Divisible par 2} & \textbf{Divisible par 3} & \textbf{Divisible par 5} & \textbf{Divisible par 10} \\
\hline
671\,890 & Oui & Non & Oui & Oui \\
\hline
666\,666\,666 & Oui & Oui & Non & Non \\
\hline
675\,985 & Non & Non & Oui & Non \\
\hline
\end{tabular}
\end{center}

\end{document}
